\documentclass{article}
\usepackage{hyperref}
\usepackage{Style}

\nocite{*} % Comentar si quiero citar
%\addbibresource{bibliografia.bib} % Quitar el comentado si quiero usar bibliografia

\begin{document}

\begin{minipage}{2.5cm}
    \includegraphics[width=2cm]{imagen_puc.jpg}
\end{minipage}
\begin{minipage}{12cm}
    {\sc Pontificia Universidad Católica de Chile\\
    Facultad de Matemáticas\\
    Departamento de Matemática\\
    Profesor: Manuel Cabezas -- Ayudante: Benjamín Mateluna}
\end{minipage}
\vspace{2ex}

{\centerline{\bf Introducción a la Geometría - MAT1304}
\centerline{\bf Ayudantía 10}}
\centerline{\bf 14 de abril de 2026}

\vspace{1cm}
\noindent\textbf{Problema 1.} Construir con regla y compás un triángulo equilatero de altura $h$
dada.

\vspace{5mm}
\noindent\textbf{Problema 2.} Dada una circunferencia con centro $O$. Sea $P$ un punto exterior
al círculo, trazar la recta tangente a la circunferencia que pasa por $P$.

\vspace{5mm}
\noindent\textbf{Problema 3.} Considere un segmento $\overline{AC}$, construya un cuadrado cuya 
diagonal sea igual a este segmento.

\vspace{5mm}
\noindent\textbf{Problema 4.} Dado un segmento $\overline{AB}$, construya el conjuntos de todos
los puntos $P$ tales que $\angle APB=45^{\circ}$.

\vspace{1cm}
\noindent\textbf{\underline{Ejercicios Propuestos:}}

\vspace{5mm}
\noindent\textbf{Extra 1.} Considere una circunferencia de centro $O$, construya un triángulo 
isósceles que lo tenga de incírculo, es decir, cuyos lados sean tangentes a la circunferecia.

\vspace{5mm}
\noindent\textbf{Extra 2.} Construir con regla y compás un triángulo $\triangle ABC$ rectángulo en 
$C$ de $m_{C}$ mediana en $C$ y lado $\overline{BC}=a$ dados.

%\printbibliography % Quitar el comentado si quiero usar bibliografia

\end{document}